In order to get used to the Machine Learning framework a first task was given to
classify EEG signals of subjects executing different tasks. There are 4
different tasks, Arithmetics, Mirror Images, Stroop, and Relax, and we expect
the brain to have different behaviour during different tasks, so, to be able to
classify the tasks from the EEG signals.

\subsubsection{The dataset}

The dataset SAM40\cite{Ghosh2022} includes brain records from 40 subjects
performing different tasks. The EEG cap is a 32-channel Epoc Flex gel kit from
the Emotiv company, which samples at 128 Hz.

\cref{fig:example-input} shows an input example, and all the EEG input in this
report will have be similar to this one (number of channels, sampling rate and
total recording time can change).

\insertfigure
    {assets/images/classification-input-example}
    {width=0.6\textwidth}
    {Sample 0 (Arithmetic task)}
    {fig:example-input}

\subsubsection{EEG classification code structure}

A big effort was made to have a clean, modular and reproducible code.  A
significant amount of time was used to build my own python ml-util module that
builds on top of PyTorch.

The designed file structure is the following:
\begin{verbatim}
. README.md
. Makefile
- dataset           (Contains datasets)
  + ...
. main.ipynb        (Notebook used for demonstration)
+ the_ml            (python scripts for models, dataset 
                     loading or preprocessing functions)
+ results           (To save results)
+ ml_util           (The custom python module)
+ checkpoints       (Stores checkpoints checkpoints)
+ config.json       (Basic hyperparameters or paths)
\end{verbatim}

The Makefile makes uses of the nbconverter command, it can generate a python
script so the whole script can be run without jupyter notebook (This was needed
because there was no port opening except the one for SSH on the supercomputer so
it was impossible to access the jupyter notebook interface).  It can also run
the notebook (without any GUI) and make a PDF out of it as a report.

This makes the training of any machine learning model as simple as the following
code.


\lstinputlisting[language=Python, caption="Basic code example showing ml\_util capabilities"]{assets/images/module_train.py}


The checkpoint function takes a callable function and a checkpoint name as
argument, if the checkpoint name file exists in the checkpoints folder then the
file is loaded and returned. Otherwise the callable function is computed and the
checkpoint file is created.

All of this clean code architecture was made to ensure reproducibility and
readability. Extending the python module to more complex tasks than
classification would require the researcher to make herited classes from the
ml\_util classes, but they would prefer making unclear code instead to keep
control of every optimization without thinking too much. This was kind of an
attempt at copying the Pytorch Lightning module but it made a good training with
PyTorch.

This file structure is inspired from the Cookiecutter Data Science template
\cite{ccds}. It serves to make the file structure common to many project so the
code is easy to read and reproducible. This is a major concern as for many
papers I had to reproduce in this internship it took a lot of time when it was
possible. One other barrier to reproduce the papers was the datasets
availability but this is for good reasons as there are privacy related reasons
that they are not publicaly available and require to make an appliance.

\subsubsection{Model for EEG classification}

\insertfigure
    {assets/images/general-classification-model}
    {width=\textwidth}
    {Overview of the classification model}
    {fig:general-classification-model}

First of all it is inconvenient to feed the raw inputs to our model, it's better
to do feature extraction before.  We features we want to extract are the alpha,
beta, gamma, delta, theta frequency bands that are very commonly used in
neuroscience, thus, we preprocess the time-frequency 5-bands spectrogram to feed
as a multi channel (one channel for each electrode) image. Indeed each sample
gets converted to a 32 channels 5x101 image, 5 corresponds to the five frequency
bands and 101 are the time intervals.

In \cref{fig:sample0-preprocessd} we show such a preprocessd sample recorded
during the arithmetic task.

\insertfigure
    {assets/images/eeg-classification-sample0-preprocessed}
    {width=\textwidth}
    {Preprocessed sample 0 (Arithmetic task)}
    {fig:sample0-preprocessd}

This multi-channels image is then fed into a custom Convolutional Neural Network
(CNN) with the architecture shown in \cref{fig:customCNN}.

This is a classical CNN architecture with a succession of 2d Convolution layers,
and a 2 layers Multi-Layer Perceptron (MLP) before a Softmax loss.

BatchNorm layers are used to improve the model learning stability and the
droupout layers are used to prevent from overfitting which occured a lot in the
experiments, over 90 percent training acccuracy was reached for 25 percent
validation accuracy.

\insertfigure
    {assets/images/custom-CNN}
    {height=0.5\textheight}
    {Custom CNN for EEG classification}
    {fig:customCNN}

\subsubsection{Training the EEG classifier}

The dataset was split, 70 percent was used for training, 15 percent for
validation, and 15 percent for test.

The validation is not trained on but serves as a test during training to tweak
the hyperparameters to detect overfitting for example.

The Adam optimizer was used, L2 regularization was added to the loss to prevent
overfitting, and every epoch, the learning rate is reduced by 1 percent.

We trained on 1000 epochs, it was enough to have a stable loss as shown in
\cref{fig:classification-hist-loss}.

\insertfigure
    {assets/images/classification_history_loss}
    {width=0.7\textwidth}
    {History of loss over the epochs}
    {fig:classification-hist-loss}

We have profiles of accuracy where we can see no clear overfittng
\cref{fig:classification-hist-acc}.

\insertfigure
    {assets/images/classification_history_acc}
    {width=0.7\textwidth}
    {History of accuracy over the epochs}
    {fig:classification-hist-acc}

\subsubsection{Results of the classifier}

The evaluation of the model on the test split of the dataset gives a 56.94
percent accuracy in the classification, the confusion matrix
\cref{fig:confusion-matrix} shows that some tasks were easier to classify.

\insertfigure
    {assets/images/classification_confusion-matrix}
    {width=0.6\textwidth}
    {Confusion matrix of our classification model on the test split}
    {fig:confusion-matrix}

We can see that the model is very good (90 percent accuracy) at classifying the
stroop effect but is pretty bad for the other tasks, especially the arithmetic
task. This can either be due to the fact that the model is not good enough or
that the EEG signals don't capture relevant enough information that can help
classify the task.